\chapter{Related Work}\label{ch:related}

This chapter situates Mol among the traditions it draws on: the classical theory of sequential machines, categorical semantics, term rewriting, and the engineering literature on high-performance dataplanes. For each field we state what it contributes, where Mol agrees with it, and where the thesis diverges. The comparison is meant to be a map, not a survey; complete bibliographies for each area are cited.

\section{Sequential Machines: Mealy and Moore}
The deterministic Mealy machine --- a state space $S$, input $A$, output $B$, and a transition-output function $S \times A \to B \times S$ --- is exactly the data of a morphism of Mol, introduced by \cite{mealy} in 1955 for the synthesis of sequential circuits; the output-decorated variant with output a function of state alone is due to \cite{moore}. Mol differs from the classical treatment in two ways. First, the morphisms of Mol are \emph{equivalence classes} of machines under the state-space bijection condition (theorem \ref{nt:2}), so state is an implementation detail and only the input--output behavior is observable. Second, Mol is organized as a category with a tensor product, so the modular structure of dataplane design --- composition and parallelism --- is first-class rather than derived. Recent work continues in the same direction: effectful Mealy machines \cite{arxiv-mealy} combine Mealy structure with a monad, which in our setting is the content of \ref{nt:5} (EffMol as the Kleisli category of the state monad).

\section{Coalgebra and Automata Theory}
Behavioral equivalence and state minimization are the core business of coalgebra. The minimal-realization theorems \ref{nt:38}--\ref{nt:41} restate, inside Mol, results whose algorithmic form is classical: \cite{nerode} for the Myhill--Nerode quotient and \cite{hopcroft} for the $n\log n$ minimization algorithm; the algebraic route via derivatives is \cite{brzozowski}. The coalgebraic literature generalizes minimization and bisimulation far beyond automata --- see \cite{arxiv-coalgebra-min} for the modern state of the art in coalgebraic partition refinement --- and Mol is a special case of that program: the step function $S \times A \to B \times S$ is an endofunctor's coalgebra once the input is fixed. What Mol adds is the category structure \emph{between} coalgebras: composition, tensor, and the enrichment of \autoref{ch:enrich} operate on machines, not just on their languages.

\section{Traced Monoidal Categories}
The traced monoidal structure of \autoref{ch:trace} is due to \cite{joyalsv}, who characterized traced monoidal categories as the universal completion of symmetric monoidal categories under feedback and established the sliding, vanishing, superposing, and yanking laws used here as \ref{nt:34}--\ref{nt:36}. Modern treatments --- see \cite{arxiv-traced-monoidal} --- reconstruct trace as a structure in profunctor categories and connect it to the geometry of interaction. The thesis's contribution is applicative: fixed-iteration loops in a dataplane are traces, the Banach contraction principle (\cite{banach}) gives their convergence, and the trace laws license the hoisting and flattening optimizations of \autoref{ch:kleisli} and \autoref{app:b}.

\section{Lawvere Theories}
Lawvere's \emph{functorial semantics of algebraic theories} \cite{lawvere} identifies an equational theory with a small category with finite products, and its models with finite-product-preserving functors. The modern correspondence between Lawvere theories and finitary monads is treated in \cite{arxiv-lawvere}. The thesis applies the framework, rather than extending it: the ring, parser, and transport theories of \autoref{ch:lawvere} are Lawvere theories, their models in Mol are the conforming components, and the completeness of the free ring model (\ref{nt:9}) is the familiar universal-algebra fact restated in Mol. The benefit of the categorical form is that models and rewrite rules are then subject to the machinery of \autoref{ch:rewrite}.

\section{Operads, Multicategories, and Clubs}
N-ary operations with typed composition are organized by operads and multicategories; the standard references are \cite{leinster} and, for the club formalism that combines operadic composition with substitution, \cite{kelly}; see also \cite{arxiv-operads} for the connection between operads, clones, and monads. Batch operations (theorem \ref{nt:46}) are precisely the operadic operations of \autoref{ch:operad}: a batch of $n$ inputs is an $n$-ary arrow, and reassociation of batching is the associativity law of operadic composition. Kelly's clubs enter at \autoref{ch:rewrite}, where typed composition trees with substitution are the objects over which the rewrite system acts.

\section{Linear Logic and Geometry of Interaction}
Girard's linear logic \cite{girard} makes resource sensitivity --- using each input exactly once --- a structural principle, and the geometry of interaction gives it an operational, token-passing semantics; the modern development of the program is \cite{arxiv-linear-logic-goi}. The thesis connects to this line in three places: cut elimination as deforestation (\ref{nt:49}), the zero-allocation guarantee (\ref{nt:50}) as the engineering reading of linearity, and the trace structure of \autoref{ch:trace}, which is the categorical shadow of feedback in the geometry of interaction. Mol's linear fragment is thus a small, decidable instance of the linear-logic programme, specialized to stateful dataflow rather than proof theory.

\section{Effectus Theory and Markov Categories}
Effectus theory \cite{arxiv-effectus} axiomatizes partiality and measurement in categorical terms, and Markov categories \cite{arxiv-markov-categories} do the same for probability; both are cousins of Mol in the family of "categories of processes." The connection points are the effectful fragment EffMol (theorems \ref{nt:5}, \ref{nt:7}) and the quantitative enrichments of \autoref{ch:enrich}, whose cost vectors are a deterministic analogue of the Markov-kernel structure: where a Markov category composes stochastic kernels, Mol composes cost-annotated deterministic transitions. The fixed-point technology differs accordingly --- Banach contraction and the Dobrushin coefficient \cite{dobroushin} for Markov updates (\ref{nt:45}) in place of probabilistic convergence.

\section{Bisimulation as Congruence}
The congruence theorem \ref{nt:6} is the thesis's use of a principle due to \cite{park}: bisimulation is the finest congruence on processes that respects observable behavior. Park's setting is nondeterministic and concurrent; Mol's is deterministic and sequential, which makes the congruence property elementary to prove (the proof of \ref{nt:6}) while keeping its practical force: equivalent machines are interchangeable in every context, which the testing strategy of \autoref{ch:testing} exploits and the minimization of \autoref{ch:coalgebra} depends on.

\section{Implicit and Descriptive Complexity}
The complexity-theoretic placement of Mol is twofold. Descriptively, finite-state molecules realize exactly the regular behaviors and minimal molecules realize minimal DFAs (\ref{nt:31}); this is the classical result of automata theory in Mol's vocabulary. Implicitly, the resource-sensitive typing of \autoref{ch:linear} belongs to the tradition of \cite{arxiv-implicit-complexity}, which characterizes complexity classes by the absence of structural operations rather than by explicit bounds; the zero-allocation theorem \ref{nt:50} is the dataplane's version of such a characterization. The accessible-category background of \autoref{ch:prelim} is \cite{adamek}.

\section{High-Performance Dataplanes}
The engineering literature on dataplanes is extensive and largely empirical: zero-copy networking, batching, SIMD packet processing, and kernel-bypass I/O are surveyed and compared in \cite{arxiv-dataplane}. The thesis's relation to this body of work is complementary rather than competitive. The empirical literature reports what works and how fast; Mol states, in theorem form, \emph{why} the known optimizations are sound --- fusion, batching, and hoisting are not heuristics but consequences of the algebra of \autoref{ch:rewrite}, \autoref{ch:operad}, and \autoref{ch:kleisli} --- and the standard that accompanies this thesis turns those theorems into enforceable policies. Where the empirical literature is the measurement apparatus, Mol is the specification the measurements certify.
EOF